\documentclass[twoside]{article}
\usepackage{package}
\usepackage[T1]{fontenc}
\usepackage{subfiles}
% Silenciar la advertencia de versalita + negrita inexistente
\WarningFilter{latex}{Some font shapes were not available}
\WarningFilter{latex}{Font shape `T1/lmr/bx/sc'}

% Silenciar la advertencia de caracteres perdidos en el nullfont (para 01_introduccion.tex)
\WarningFilter{latex}{Missing character: There is no}
%Nuevos comandos
\renewcommand{\thefootnote}{\arabic{footnote}}
\raggedbottom



\begin{document}
\begin{center}
    \thispagestyle{empty}
    \large{\textsc{Procesamiento de Datos y Modelos Avanzados en Simulink}}\\[0.5cm] 
    \large{Fabrizio Benavidez\footnote[1]{\texttt{fabrizio.benavidez@mi.unc.edu.ar}}}\\[0.5cm]

     \href{https://github.com/Fabopy/Matlab-curso-}{%
        {\Large\faGithub}\quad
        \nolinkurl{https://github.com/Fabopy/Matlab-curso-}%
    }\\[0.5cm]

    \large{\footnotemark[1]\textit{Facultad de Ciencias Exactas, Físicas y Naturales, Ingeniería Aeronáutica}}
    \\[0.5cm]

    \begin{minipage}{0.7\textwidth}
    \hrule
    \vspace{0.2cm}
        \normalsize{\textsc{Resumen} El siguiente trabajo consiste en el análisis de datos e imágenes con MATLAB. En el mismo se presenta el análisis de datos de un hoja de cálculo (Excel), así como el de una figura jpg. 
        
        Además, también se analiza otros modelos de Simulink un poco más avanzados. Se ven los bloques \textit{From Workspace}, \textit{LTI System} y \textit{Fcn}}\\[0.2cm]
        \normalsize{\textit{Palabras claves}: Datos, imágenes, MATLAB, Simulink, bloques \ }
    \vspace{0.2cm}
    \hrule
    \end{minipage}
    
    \vfil

    \includegraphics[width = 0.4\textwidth]{imagenes/Matlab_Logo.png}\\[0.5cm]

    \vfil

    \today\\[0.2cm]
    \large{Semana 9}
\end{center}







\newpage

\section*{Guía 17: Adquisición y procesamiento de datos} 
\subfile{secciones/guia17}


\newpage

\section*{Guía 18: Modelos avanzados de Simulink} 
\subfile{secciones/guia18}














\newpage
% \thispagestyle{empty}
% \bibliography{biblioteca.bib}
% \bibliographystyle{plain}


\end{document}